\section{Introduction}

The UK's housing sector is at a critical crossroads, facing a multitude of challenges that are deeply intertwined with the nation's broader socio-economic and environmental goals. Among these challenges, the need to decarbonize the housing stock to meet Net Zero commitments by 2050 stands out as a pivotal issue. A significant portion of the UK’s residential buildings, particularly those in the rental sector, are currently underperforming in terms of energy efficiency since as of 2022, homes with an Energy Performance Certificate (EPC) rating above Band C only represented 42\% and 37\% of the total stock in England and Wales, respectively \cite{office_for_national_statistics_age_2022}. This underscores the vast potential—and necessity—for improvements in the energy performance of the UK’s housing stock.\\

In the 2024 elections’ manifesto, the Labour Party mentioned its intentions to invest an additional £6.6 billion over the current parliament to increase home energy efficiency as part of its 'Energy Independence Act' \cite{noauthor_make_nodate} as part of their plan to take increasing actions to reduce energy poverty and increase UK home energy efficiencies in their pursuit of achieving NZ \cite{scott_kings_2024}. This follows decades of multiple policy agendas that have been unsuccessful at engaging privately rented dwellings to improve their energy efficiency \cite{hope_attitudes_2014}. In October 2021, the UK Government released the Heat and Buildings Strategy, outlining its approach to decarbonizing homes and buildings \cite{department_for_energy_security_and_net_zero_desnz_and_department_for_business_energy_and_industrial_strategy_beis_heat_2021}. A primary objective of this strategy is to enhance the energy efficiency of as many homes as possible across England and Wales, with the goal of achieving an EPC rating of Band C by 2035, provided this is practical, cost-effective, and financially viable. Additionally, the strategy aims for as many fuel-poor homes as reasonably practicable to reach an EPC Band C by 2030. While there are existing minimum energy efficiency standards (MEES) for privately and social-rented homes in the UK and Wales \cite{rowe_housing_2024}, new and existing private tenancies were required to be certified as Band E or higher since April 2018 and 2020, respectively (The Energy Efficiency (Private Rented Property) (England and Wales) Regulations, 2015). Should landlords fail to achieve MEES, they may be required to pay a fine up to £5,000 \cite{department_for_energy_security_and_net_zero_2023_2023}. While the UK government does offer grants to partly cover the costs of retrofitting a dwelling up to standards, such as the Boiler Upgrade scheme which subsidises the deployment of heat pumps, and no property owner is currently obligated to spend in excess of £3,500 on reaching the MEES \cite{department_for_energy_security_and_net_zero_2023_2023}. In other words, when a property is unable to reach the MEES within the given expenses, landlors must only make improvements up to that amount \cite{department_for_energy_security_and_net_zero_2023_2023}.   In September 2020 and January 2021, the government consulted whether to amend the current MEES to Band C by 2025 for new leases and 2028 for existing ones and revise upward the £3,500 threshold \cite{department_for_energy_security_and_net_zero_desnz_and_department_for_business_energy_and_industrial_strategy_beis_heat_2021}. In September 2023, then Prime Minister Rishi Sunak announced that the target for UK rented homes would be scratched,  which was later brought back by the Labour Party \cite{noauthor_make_nodate}.\\


Retrofitting homes has been referred to as the most significant challenge to meet the UK Net Zero goals  \cite{alabid_review_2022, uk_green_building_council_ukgbc_retrofit_2021}. As with all infrastructural transformations, sustainable structural retrofits are difficult to implement at a rapid and large scale due to structural, technical and behavioural path dependences \cite{dauda_understanding_2022}. Committing to green retrofits to limit energy consumption and resulting carbon emissions has been not only historically difficult \cite{lowe_technical_2007, chaves_building_2017}, but also complicated due in part to UK property law which acts as an additional barrier for dwellings with mixed tenures and leaseholders (particularly flats) as it requires consent from all stakeholders involved \cite{weatherall_property_2018, bright_exploring_2019}. This is further exacerbated by properties' long lifetime and slow retrofitting rates, making it difficult and costly for the housing sector to shift from its high-emitting paradigm, as it is estimated that nearly 80\% of properties were built to standard levels below those required to meet  Net Zero targets \cite{warren_can_2024, uk_green_building_council_ukgbc_retrofit_2021}. \\

While incentives to commit to expensive retrofits are mixed, there is evidence that mandatory energy efficiency disclosure have resulted in green premiums, increased energy investments, and helped correct some marker failures in the housing sector \cite{cassidy_how_2023, myers_mandatory_2022}. Green Premiums in particular have been extensively studied as consumer’s willingness-to-pay for homes with higher energy efficiency at the time of purchase \cite{harrison_hedonic_1978, bisello_measuring_2020, dalton_green_2018, ghosh_homebuyers_2024, fuerst_is_2020, kahn_capitalization_2014} or rent \cite{bruegge_does_2016, galvin_rental_2023, sieger_inefficient_2023}. While most empirical results concur green premiums on rent and sales prices to be positive (usually around 6\% and 7.6\%, respectively), there are mixed interpretations about the cost-effectiveness of undergoing energy efficiency improvements for landlords \cite{dalton_green_2018}. This is especially the case for deep retrofitting costs which may have long and unpredictable pay-back periods, at the time of resale. Nearly 28\% of all property owners refusing to retrofit do not do so because of the high costs of retrofitting \cite{office_for_national_statistics_age_2022}. It is important to note that estimating these green premiums suffers from confounding factors such as increased demand for homes outpace the supply of newly built homes in the UK \cite{aldridge_housing_2018, nazir_comparison_2020}, and this widening supply gap exacerbates housing shortages, contributing to rising property prices \cite{green_policies_2017}. COVID-19 has also exacerbated the housing crisis, as demand for bigger properties which could accommodate for remote working increased, even well after the pandemic \cite{schwartz_covid-19s_2023}. The resulting housing instability and inaccessibility have far-reaching implications, affecting socio-economic, political, and health systems \cite{murphy_housing_2024, bone_neoliberal_2014, banks_house_2010, mckee_generation_2020}. A notable consequence of this crisis is the decline in housing mobility, particularly among younger generations, which has led to a sharp increase in inequality and reduced disposable income \cite{howard_seven_2024}. This trend is expected to persist, further deepening disparities in the housing market and willingness to commit to green retrofits \cite{green_policies_2017, sabater_age_2023}. \\

In parallel with the housing supply crisis, the private renting sector has expanded rapidly over the past two decades \cite{field_reappraising_2014, bailey_poverty_2020}. As of 2022, approximately 4.6 million households, or 19\% of tenures, were in the private rental sector in England \cite{department_for_levelling_up_housing_and_communities_dluhc_english_2022}. Despite the recent trends in unaffordability, there is a still an ever-increasing demand for properties as people continue to want to own their own homes \cite{aldridge_housing_2018}. A significant proportion of these private renters—over 62\%, equating to around 2.8 million households—expressed plans to buy a home in the future, a notable increase from 56\% in 2019/20 \cite{department_for_levelling_up_housing_and_communities_dluhc_english_2022}. Private renters often face additional challenges, including higher costs for utilities as over half of private renters pay their gas bills via direct debit, and two-thirds do the same for electricity (EHS, 2022). The most common type of tenancy in the domestic sector is a short-hold tenancy (EHS, 2022) leading to increased insecurity amongst the younger and lower-income groups \cite{mckee_generation_2020}. Moreover, private renters generally live in homes with poorer energy performance compared to social renters, further compounding their financial burdens and quality of living (EHS, 2022). \\


Overall, the residential sector is a significant contributor to the UK's energy consumption and carbon emissions since UK households consume on average 3,500 kWh of electricity and 12,000 kWh of gas annually \cite{department_for_business_energy__industrial_strategy_dbeis_energy_2022}. A substantial portion of these emissions stems from heating and electricity use within homes, which are primarily driven by fossil fuels. Notably, space heating accounts for about 63\% of household energy consumption, while water heating, lighting, and appliances collectively contribute to the remaining 37\% \cite{palmer_united_2014}. Energy prices have fluctuated considerably in recent years, with significant increases observed during periods of geopolitical instability, such as the Russian invasion of Ukraine which exacerbated energy price volatility as Europe faces simultaneous disruptions in gas supply and a shift towards alternative energy sources, including renewables \textbf{(Drax Group, 2023)}. The penetration of renewables in the UK grid has been a crucial factor in buffering these price shocks, with renewables accounting for 40\% of electricity generation in 2022, up from 29\% in 2017 \textbf{(National Grid ESO, 2023)}. Retrofitting strategies have shifted from energy conservation measures (i.e., improving wall insulation and double-glazed windows) towards integrating smart technology measures to enhance energy efficiency and integration of renewable technologies \cite{ejidike_global_2023}. Space heating and hot water are responsible for most of the energy consumption in UK households and have hence been considered promising avenues for end-user decarbonization \cite{lowe_technical_2007, kesicki_costs_2012}. Despite being a hard-to-abate sector, reducing end-use emissions associated with heating and powering the domestic building stock seems the most cost-effective measure \cite{kesicki_costs_2012, royapoor_towards_2023}. Technical predictive scenarios have suggested the necessary CO2 emission reductions could be achieved through technology readily available when accompanied by the decarbonization of the system-level electric grid and increasing switch away from fossil fuels \cite{lowe_technical_2007, kesicki_costs_2012, miu_simple_2018}. However, the current retrofitting rates are far below the expected levels to reach NZ targets. Hence, there is an increasing need for policies to enable technological uptake within domestic dwellings if the UK is to meet its net zero targets \cite{wang_innovative_2022, miu_simple_2018}. \\

Heightened concerns regarding the net effect of climate change on domestic energy usage has also led to increased interest in energy conservation and preventative measures to mitigate the anticipated rise in insulation and air-conditioning needs \cite{de_cian_households_2019, poritt_scopus_2012}. In these scenarios, financial motivations, particularly the payback period and projected energy savings, often dominate the decision to renovate, suggesting a "double-win" strategy that reduces both energy costs and carbon emissions \cite{joint_research_centre________________________________________european_commission_building_2020}. The effectiveness of energy efficiency labels to affect property prices has been shown to vary by climatic conditions \cite{taltavull_de_la_paz_green_2019} and age of target audience \cite{wilhelmsson_how_2023}, further adding to the heterogeneous and unpredictable returns from the perspective of property owners. The cognitive perceptions and financial situation of homeowners may significantly impact what is considered a cost-effective energy renovation \cite{galassi_role_2017, ud_din_assessing_2023, wilson_why_2015}. Overall, for homeowners to commit to green retrofits, their preferences and values must align with their perceived financial returns on investment, in that they must be able to expect either an increase in the future sale price and/or a reduction in operational costs (i.e., energy bills) over a sufficiently short pay-back period \cite{wilson_why_2015, wang_holistic_2022}, which results in a demographically and geographically heterogenous decision-making landscape \cite{pelenur_motivations_2014}. \\ 


This study contributes to the ongoing discourse on housing policy and energy efficiency by providing a nuanced understanding of how imposing EPC standards might influence the UK housing market. This research offers insights into how such regulations could not only help achieve environmental targets but also address housing supply and affordability issues. Through detailed empirical analysis, this study aims to inform future policy decisions, ensuring that the UK can meet its Net Zero goals while also fostering a more equitable and efficient housing market. \\

The contributions to the literature are threefold. First, this study provides an empirical estimation of the costs of retrofitting and their associated impact on energy consumption, carbon emissions and overall energy rating. Second, this study attempts to quantify the green premium associated with higher energy ratings as stated in their Energy Performance Certificate (EPC) at point of sale using transaction data. Third, the study estimates the potential effect of an increase on the supply side on property prices using listing data. \\

The remainder of this paper is structured as follows. The next section describes the different datasets used in this study. The third section presents the methodology alongside the main results. Section 4 concludes. 


